\section{技术选型与协作开发}
\begin{frame}{开发技术栈与选型理由}
\small
\begin{tabularx}{\textwidth}{lXX}
\toprule
层次 & 技术与工具 & 选择理由 \\
\midrule
前端 & Vue 3、TypeScript、Element Plus & 支持中文表单、运行状态和证据浏览，类型约束便于联调。\\
服务与数据 & Python、FastAPI、Pydantic、SQLite & 共用接口模型，支持事件推送，本地部署依赖较少。\\
智能体与浏览器 & LangGraph、HTTPX、Playwright & 显式管理状态与预算，可获取页面观察、截图及执行轨迹。\\
环境管理 & Docker Compose、uv、pnpm & 隔离被测应用，锁定依赖，提高团队环境一致性。\\
质量与文档 & pytest、Ruff、前端测试、XeLaTeX & 分层验证执行逻辑、界面流程，维护可版本化报告。\\
\bottomrule
\end{tabularx}
\end{frame}

\begin{frame}{Git 协作与开发过程留痕}
\begin{itemize}
\item \textbf{协作平台：}计划使用 GitHub，按需求建立 Issue，记录负责人、验收条件和关联里程碑。
\item \textbf{开发流程：}从主分支建立短分支，经 Pull Request 评审并通过检查后合并。
\item \textbf{过程数据：}保留提交、讨论、代码评审与 CI 日志，使需求到交付的过程可查。
\item \textbf{版本管理：}提交依赖锁文件、被测版本清单与报告源码，实验结果注明代码版本。
\item \textbf{质量控制：}使用 Conventional Commits 和 Git hooks，统一执行检查脚本。
\end{itemize}
\vfill
\small 当前仓库已提供协作约定和 CI 配置；远端仓库、分支保护及团队访问需在协作启动时完成。
\end{frame}
